\chapter{Extension: when merging two cohorts pays off}\label{ch:merging}
% Not in the thesis: Result A of notes/merging-cohorts.md (Case A).
% Lean: Credibility/Merging.lean

Not in the thesis; it answers the question of
Remark~\ref{rem:equalZ_grouping}. Take two cohorts on their own, as in
Chapter~\ref{ch:two}, with $\mu$ estimated from the pair. Let
$v_j = \sigma_j^2$ be the error variances, $c$ the error covariance
($c = \sigma_1\sigma_2\rho$ in Chapter~\ref{ch:two}),
$V = \operatorname{Var}(\hmu_1 - \hmu_2)$ and $d = \mathbb{E}(\mu_1-\mu_2)^2$
the true gap between the levels. An unbiased linear predictor of $\mu_1$ is
$\hmu_1 - \beta(\hmu_1 - \hmu_2)$. A working model with gap variance $x$ uses
$\beta(x) = (v_1 - c)/(V + x)$: keeping the cohorts separate (the thesis model)
is $x = 2\tau^2$, merging them ($\mu_1 = \mu_2$) is $x = 0$. Under the thesis
model $\beta(2\tau^2)$ is the weight $(1 - z_1)b_2$ of \eqth{32}.

\begin{definition}[Merging quantities]\label{def:merging}
  \lean{Credibility.Merging.V, Credibility.Merging.moments, Credibility.Merging.mse,
        Credibility.Merging.weight, Credibility.Merging.Q,
        Credibility.Merging.excessLoss}
  \leanok
  \uses{def:bestZ}
  $V = v_1 + v_2 - 2c$ and $Q = (v_1 - c)^2 + (v_2 - c)^2$. The predictor
  $\hmu_1 - \beta(\hmu_1 - \hmu_2)$ is the blend $z\hmu_1 + (1-z)\hmu_2$ with
  $z = 1 - \beta$, so its mean squared error $\operatorname{mse}_1(\beta)$ is
  \eqth{6} at $z = 1 - \beta$, for the second moments
  $\mathbb E(\hmu_1 - \mu_1)^2 = v_1$,
  $\mathbb E(\hmu_2 - \mu_1)^2 = v_2 + d$ and
  $\mathbb E(\hmu_1 - \mu_1)(\hmu_2 - \mu_1) = c$; likewise for cohort 2. The
  weight under a working gap variance $x$ is $\beta_j(x) = 1 - z^*$, with
  $z^*$ the factor \eqth{7} when the gap is $x$. The excess loss of the
  working value $d_w$ is
  $\Delta L(d_w) = \sum_{j=1}^2 \bigl(\operatorname{mse}_j(\beta_j(d_w)) -
  \operatorname{mse}_j(\beta_j(d))\bigr)$.
\end{definition}

\begin{lemma}[Closed forms]\label{lem:merging_apply}
  \lean{Credibility.Merging.mse_apply, Credibility.Merging.weight_eq}
  \leanok
  \uses{def:merging}
  $\operatorname{mse}_j(\beta) = v_j - 2\beta(v_j - c) + \beta^2(V + d)$, and
  $\beta_j(x) = (v_j - c)/(V + x)$ if $V + x \ne 0$.
\end{lemma}
\begin{proof}
  \leanok
  Expand \eqth{6} and \eqth{7}; the curvature is $V + x$.
\end{proof}

\begin{lemma}[Excess-loss identity]\label{lem:mse_eq}
  \lean{Credibility.Merging.mse_eq}
  \leanok
  \uses{def:merging}
  If $V + d \ne 0$, then
  $\operatorname{mse}_j(\beta) = \operatorname{mse}_j(\beta^*) + (V + d)(\beta - \beta^*)^2$
  with $\beta^* = (v_j - c)/(V + d)$.
\end{lemma}
\begin{proof}
  \leanok
  \uses{lem:errVar_eq}
  Lemma~\ref{lem:errVar_eq}: \eqth{6} is a parabola with vertex at \eqth{7}
  and curvature $V + d$.
\end{proof}

\begin{lemma}[Closed form of the excess loss]\label{lem:excessLoss_eq}
  \lean{Credibility.Merging.excessLoss_eq}
  \leanok
  \uses{lem:mse_eq}
  If $V + d \ne 0$ and $V + d_w \ne 0$, then
  \[
    \Delta L(d_w) = Q\,(V + d)\Bigl(\frac{1}{V + d_w} - \frac{1}{V + d}\Bigr)^2.
  \]
\end{lemma}
\begin{proof}
  \leanok
  \uses{lem:mse_eq, lem:merging_apply}
  By Lemma~\ref{lem:mse_eq} cohort $j$ contributes $(V + d)(\beta_j(d_w) - \beta_j(d))^2$,
  and $\beta_j(d_w) - \beta_j(d) = (v_j - c)\bigl(1/(V + d_w) - 1/(V + d)\bigr)$.
\end{proof}

\begin{theorem}[Result A: midpoint rule]\label{thm:merge_iff}
  \lean{Credibility.Merging.merge_iff}
  \leanok
  \uses{lem:excessLoss_eq}
  Let $\tau^2 > 0$, $V > 0$, $Q > 0$ and $d \ge 0$. Merging beats keeping the
  cohorts separate, $\Delta L(0) < \Delta L(2\tau^2)$, if and only if
  \[
    d\,(V + \tau^2) < V\tau^2, \qquad\text{i.e.}\qquad d < \frac{V\tau^2}{V + \tau^2}.
  \]
\end{theorem}
\begin{proof}
  \leanok
  \uses{lem:excessLoss_eq}
  The common factor $Q(V + d) > 0$ cancels. With $t = 1/(V + d)$,
  \[
    \Bigl(\frac{1}{V + 2\tau^2} - t\Bigr)^2 - \Bigl(\frac1V - t\Bigr)^2
    = \frac{4\tau^2\bigl(V\tau^2 - d(V + \tau^2)\bigr)}{V^2(V + 2\tau^2)^2(V + d)}.
  \]
  So merging wins iff $t$ lies above the midpoint of $1/V$ and $1/(V + 2\tau^2)$.
\end{proof}

\begin{remark}\label{rem:merge_reading}
  \uses{thm:merge_iff}
  The condition involves only $V$, $\tau^2$ and $d$, not $z_1$ against $z_2$.
  Correlation enters through $V = \sigma_1^2 + \sigma_2^2 - 2\sigma_1\sigma_2\rho$:
  $\rho > 0$ makes merging harder. Under the thesis model $d = 2\tau^2$, and
  $2\tau^2(V + \tau^2) > V\tau^2$, so merging never pays. Because $Q$ is a
  common factor, each cohort on its own prefers the same choice.
\end{remark}

% ---------------------------------------------------------------------------
